\section{Meshing and Convergence Analysis}

\subsection{Mesh Description}

% Type of mesh: automatic, tetrahedral, hexahedral...
% Local refinements, element size, element type

% TODO: Include screenshot of the global mesh
% \begin{figure}[H]
%     \centering
%     \includegraphics[width=0.8\textwidth]{mesh_global}
%     \caption{Global mesh of the geometry.}
%     \label{fig:mesh_global}
% \end{figure}

% TODO: Include screenshot of local refinements if any
% \begin{figure}[H]
%     \centering
%     \includegraphics[width=0.8\textwidth]{mesh_refined}
%     \caption{Detail of the refined mesh zones.}
%     \label{fig:mesh_refined}
% \end{figure}

% State element type, characteristic size, approximate number of nodes and elements

\subsection{Convergence Analysis}

Before accepting the modal analysis results, a mesh convergence study has been carried out. The modal analysis has been repeated with several meshes of different sizes and the natural frequency of one vibration mode has been compared. Further refinement beyond an element size of 0.003125\,m was not possible due to the node and element count limitations of the ANSYS Student licence.

An acceptable convergence criterion is a relative variation of less than 5\% between two consecutive refinements:
\begin{equation}
    \Delta f_i = \frac{|f_{i,\text{refined}} - f_{i,\text{previous}}|}{f_{i,\text{refined}}} \times 100
    \label{eq:convergence}
\end{equation}

The complete convergence results are presented in \Cref{tab:convergence}. All meshes satisfy the 5\% criterion from the second refinement onward, and the relative error between consecutive refinements falls below 0.1\% for element sizes of 0.015\,m and smaller, confirming that the solution has converged.

\begin{table}[H]
\centering
\caption{Mesh convergence analysis for the selected mode.}
\label{tab:convergence}
\begin{tabular}{cccc}
\toprule
Element size [\si{\metre}] & $f_i$ [\si{\hertz}] & Variation [\%] & Total error [\%] \\
\midrule
0.500000 & 101.3230 & ---   & 3.9977 \\
0.100000 &  98.7641 & 2.59  & 1.3713 \\
0.050000 &  98.6038 & 0.16  & 1.2067 \\
0.025000 &  97.6369 & 0.99  & 0.2143 \\
0.022000 &  97.5303 & 0.11  & 0.1049 \\
0.020000 &  97.5306 & 0.00  & 0.1052 \\
0.015000 &  97.5226 & 0.01  & 0.0970 \\
0.012500 &  97.4837 & 0.04  & 0.0570 \\
0.006250 &  97.4384 & 0.05  & 0.0106 \\
0.003125 &  97.4281 & 0.01  & 0.0000 \\
\bottomrule
\end{tabular}
\end{table}

\Cref{fig:convergence_frequency} shows the natural frequency as a function of mesh refinement. As the mesh becomes finer, the frequency converges to a stable value around 97.4\,Hz.

\begin{figure}[H]
    \centering
    \includesvg[width=0.85\textwidth]{frequency}
    \caption{Natural frequency vs.\ inverse mesh size.}
    \label{fig:convergence_frequency}
\end{figure}

\Cref{fig:convergence_error_measures} presents the relative error between consecutive refinements. After the initial coarse meshes, the error drops below 1\% for element sizes of 0.025\,m and smaller.

\begin{figure}[H]
    \centering
    \includesvg[width=0.85\textwidth]{error_measures}
    \caption{Relative error between consecutive refinements vs.\ inverse mesh size.}
    \label{fig:convergence_error_measures}
\end{figure}

\Cref{fig:convergence_error_total} shows the total relative error with respect to the finest mesh. The error falls below 0.1\% for element sizes of 0.022\,m and smaller, indicating that the solution is mesh-independent at that level of refinement.

Based on the convergence study, the mesh with an element size of 0.022\,m ($1/h \approx 45.5$\,m$^{-1}$) has been selected for the subsequent analyses. As can be observed in \Cref{fig:convergence_frequency} and \Cref{fig:convergence_error_total}, the natural frequency stabilises at this point and the total error falls below 0.11\%, which is sufficiently accurate for the purposes of this study. Choosing this mesh instead of a finer one significantly reduces computational cost and time while remaining within the convergence region.

\begin{figure}[H]
    \centering
    \includesvg[width=0.85\textwidth]{error_total}
    \caption{Total relative error vs.\ inverse mesh size.}
    \label{fig:convergence_error_total}
\end{figure}